%% -----------------------------------------------------------
\section{Research Gaps and Future Directions}
%% -----------------------------------------------------------

Five critical gaps require urgent attention.

\textit{Long-term outcome tracking} is nearly absent. Smith et al.~\cite{smith_2023} establish that the field lacks evidence about sustained benefits, delayed harms, or consequences when users discontinue chatbot use. Without longitudinal data, clinical claims cannot be evaluated against real-world outcomes.

\textit{Intersectionality} receives minimal attention. Studies examine single identity categories rather than multiply-marginalized individuals who face compounded barriers. \gls{lgbtq} individuals who are also disabled, non-Western, and low-income, or neurodivergent individuals who are also racially minoritized, represent the most underserved populations, yet they appear in almost no research~\cite{liu_2023,timmons_2023}.

\textit{Crisis performance} is dangerously understudied. Limited systematic evidence exists about AI performance during acute suicidal emergencies, precisely the context where failure proves fatal~\cite{de_freitas_2024}. Existing evidence from companion AI and general chatbot studies is concerning; purpose-built crisis response has received almost no controlled investigation~\cite{smith_2023}.

\textit{Neurodivergent-specific design} remains nascent. Despite compelling evidence of unmet need and specific design requirements, the field lacks \gls{rct}s of purpose-built AI mental health tools for autistic or \gls{adhd} populations and lacks standardized evaluation protocols for neurotypical bias~\cite{killian_2023,papadopoulos_2025,dehghani_2024}.

\textit{Comparative effectiveness} research comparing AI with appropriately matched human psychotherapy for marginalized populations does not exist at the scale necessary to make evidence-based deployment recommendations~\cite{abd_alrazaq_2020}.
